% Day 08 例题 3：光滑斜面与粗糙水平面的多过程
\begin{tikzpicture}[scale=0.83,>=Latex,line cap=round,line join=round]
  \coordinate (A) at (0,2.75);
  \coordinate (B) at (3.25,0);
  \coordinate (C) at (7.35,0);

  \draw[very thick] (A) -- (B) -- (C);
  \foreach \x in {3.55,3.95,...,7.05}
    \draw[gray!65] (\x,0) -- ++(-0.13,-0.13);

  \draw[densely dashed,gray] (0,0) -- (A);
  \draw[<->,thick] (-0.38,0) -- (-0.38,2.75)
    node[midway,left,fill=white,inner sep=1pt] {$h=5\ \mathrm{m}$};
  \node[rotate=-40.2,above=4pt] at (1.53,1.29) {光滑斜面};
  \node[above=4pt] at (5.55,0) {粗糙水平面};

  \begin{scope}[shift={(0.70,2.16)},rotate=-40.2]
    \draw[fill=gray!12,thick] (-0.38,0) rectangle (0.38,0.46);
  \end{scope}
  \node[above left=1pt] at (0.58,2.68) {$v_0=0$};

  \draw[fill=gray!12,thick] (4.50,0) rectangle (5.30,0.52);
  \node at (4.90,0.26) {$m$};
  \draw[->,thick,blue] (5.30,0.26) -- ++(0.95,0)
    node[right=2pt,fill=white,inner sep=1pt] {$v$};
  \draw[->,thick,red] (4.50,0.26) -- ++(-0.85,0)
    node[left=2pt,fill=white,inner sep=1pt] {$f_k$};

  \draw[->,orange!85!black,thick,dashed] (1.18,1.74) -- (2.48,0.65);
  \draw[<->,thick] (3.30,-0.45) -- (7.00,-0.45)
    node[midway,below=2pt,fill=white,inner sep=1pt] {$x$};
  \fill (7.00,0.08) circle (1.5pt);
  \node[above right=1pt] at (7.00,0.08) {停止};
\end{tikzpicture}
